\section{Introduction}

The preceding provenance and activation papers separated three objects that
are easily conflated: a declared parent incidence, the literal terminal atoms,
and the energy that is actually generated by a fixed row coefficient
\citep{WangTPC63,WangTPC64}.  That separation resolves which denominator is
being charged.  It does not show that a coherent residual or native
representation preserves the activated energy.  The latter is a matrix
question on the range populated by the physical coefficients.

This paper closes that interface at the exact finite-dimensional level.  Its
first point is a restriction principle: the free ambient terminal atom space
is not the represented reference.  The correct carrier is the actual analysis
range, or equivalently the original row space equipped with the semidefinite
terminal form.  On this carrier the represented map has a concrete column
Gram.  Large kernels of an ambient grouping map can be irrelevant when they
miss the populated range; conversely, every individual row can be nonzero while
two active row columns cancel coherently.

The general shorted-operator formula is not new here.  It was proved and
audited in the exact closure calculus of \citet{WangTPC62}, building on the
classical shorted-operator viewpoint \citep{AndersonTrapp1975}.  We use that
formula only to distinguish two cases.  If the represented map descends through
terminalization, its terminal-null directions vanish and the short reduces to
an ordinary Gram on the actual range.  If it does not descend, deleting those
directions would be an invalid compression; the full short must be retained.

The new content is fourfold.
\begin{enumerate}[label=(\roman*)]
\item We derive the literal residual-fiber row Gram and an exact criterion for
represented positivity on the activated range.
\item We prove invariance of this Gram under every terminal-faithful pair
repair.  Repair-only parent atoms affect activation relative to the parent but
cannot be charged again after terminalization.
\item We establish a multiblock kernel-fusion bound.  Singular dyadic block
Grams can combine to a positive global Gram when their blind directions are
complementary.
\item We prove a sharp spectral-capacity law and an equally sharp non-transfer
theorem: the row-capacity certificate from the activation layer survives only
after a reducing/basis-transport statement connects the row basis to the good
spectral subspace.
\end{enumerate}

These conclusions advance the interface but not the arithmetic exponent.  In
particular, injectivity at each finite \(X\) need not give a polynomial lower
bound, and a labelled isometry need not survive physical label erasure.  The
next step must therefore study the literal collision geometry of the
residual-fiber columns, rather than infer it from activation counts.

\subsection*{Claim level and quantifiers}
The linear-algebraic results are unconditional L0 statements for the declared
finite spaces and maps.  Their fixed-\(h_0\) dictionary is L1 only when the same
literal coefficients, actual support, block labels, and raw atomic norm from
the activation paper are retained.  No result below proves a favorable
activation census, a polynomial lower singular value for the physical
unlabelled map, a parity-breaking estimate, or a prime-pair lower bound.
